\chapter{RQ2 --- Calibrated Soft-Routing}
\label{ch:rq2}

\emph{From an oracle ceiling to a realizable, auditable address.}

\section{The question RQ2 must answer}

The representation chapter (RQ1) establishes that a type-conditional spatial address --- the
position-slot $(i, M)$ for a filler, the connectivity fingerprint for a wall --- collapses the
visual-confusable set almost completely \emph{when supplied at an oracle level}: the retrieved pool
(median 76, ground truth in-pool 100\%) reorders from a realized Top-1 of 6.7\% to \textbf{78.5\%}
(fillers to 91\%), with the ground truth in the pool throughout. That ceiling is a statement about
\emph{information}: the address, if known, suffices. It says nothing about whether the address can be
\textbf{recovered} from a site photograph and a plan, nor how a \emph{noisily} recovered address
should be used. RQ2 is the gap between the two:

\begin{quote}
\emph{How much of the oracle ceiling is realizable once the address is extracted by a real, imperfect
detector --- and how is a partly-reliable address consumed without destroying the recall that the pool
guarantees?}
\end{quote}

The honest baseline is stark. The thesis reranker (G8) emits the position field as free text and, on
the held-out fillers, recovers a usable slot in \textbf{0 of 35} cases. With no slot signal the
pool-ranking falls back to storey and class, a Top-1 of \textbf{2.4\%}; even a modal-prior slot guess
--- the most generous floor --- reaches only \textbf{6.6\%} (G8's overall held-out Top-1 is 6.7\%).
The oracle says 91; the system delivers single digits. RQ2 is the work of closing that.

\section{The recall constraint: why the address must be a soft prior, not a hard filter}

The intuitive way to use an extracted address is to \emph{filter} the pool by it --- keep only
candidates whose slot matches. This is exactly wrong, and the measurement says why. Treating each
extracted field as a hard constraint multiplies the per-field recall: with measured single-field
reliabilities (storey $\approx 0.66$, class $\approx 0.50$), the joint recall of a hard conjunction
\emph{over the measured attribute fields} collapses toward $\prod r \approx 0.009$; the
position-slot's own extraction reliability is unmeasured, but its realized 0-of-35 recovery makes it
no exception. The filter evicts the ground truth from the pool far more often than it removes a
distractor, and the 100\%-recall guarantee that made the task tractable is gone. Hard determinism and
recall are in direct opposition.

The resolution --- and the first refinement RQ2 forces on its own original ``deterministic
retrieval'' framing --- is to relocate determinism. \textbf{The executor is deterministic given the
structured record; the \emph{ranking} is a calibrated soft prior inside a recall-fixed pool.} The
address never removes a candidate; it re-weights one. Recall stays at 100\% by construction, and the
unreliable signal is spent only on \emph{ordering}. This is the mechanism the rest of the section
makes precise and measures.

\section{The mechanism: a per-field confidence contract, calibrated, routed}

Every extracted field is emitted as a uniform record \texttt{\{value, confidence, source\}} (the
contract invariant; it closes the two extraction-confidence gaps noted in the system's own
limitations --- confidence underused, and no per-model calibration). The position-slot detector, a
deterministic OpenCV specialist that color-segments the plan's openings and orders them along the
host wall, produces \texttt{value} $= (i, M)$, a \texttt{confidence}, and \texttt{source = opencv}. A
routing policy assigns each field a role and a weight; the recovered slot enters the rerank as a
confidence-weighted prior, and --- crucially --- a low-confidence extraction can route to
\textbf{defer} (``I am not sure; here are the candidates'') rather than to a confident wrong answer.
The pipeline is summarized in Figure~\ref{fig:pipeline}, with the confidence-routing path traced from
the detector's \texttt{confidence} to the answer/defer gate; that path \emph{is} the
determinism$\leftrightarrow$adaptivity mechanism.

We report four results.

\paragraph{(1) One realizable extractor closes most of the oracle gap.} The position-slot detector,
scored against the image-recoverable convention (below), recovers the opening count exactly in
\textbf{83\%} of cases (29/35) and the full slot in \textbf{74\%} (joint $(i, M)$). Fed into the soft
rerank, it lifts filler Top-1 from the \textbf{6.6\%} modal-prior floor to \textbf{67.6\%}, and Top-10
to \textbf{80.9\%}, at zero recall cost --- the floor-to-realized arc the oracle gap predicted, now
delivered by a real detector rather than an oracle. The \emph{count} is reliable; what residual error
remains is concentrated in the ordering index, not in the perception of the openings.

\paragraph{Why a deterministic detector rather than the fine-tuned VLM.} A per-field re-evaluation of
the LoRA-fine-tuned multimodal extractor (G8) over the held-out set is decisive (Figure~\ref{fig:vlm}):
it recovers the \textbf{coarse prefix at 100\%} (storey and \texttt{ifc\_class}, both ${\sim}100\%$
correct) but the \textbf{discriminating structured fields at 0\%} --- the position-slot and the size
band --- with relation \emph{direction} only partial (57\%). The fine-tuned network is thus saturated
on exactly the part of the address that, by RQ1, does \emph{not} discriminate, and absent on the part
that does. The lesson is architectural, not ``fine-tune harder'': the slot and size are
\textbf{delegated to deterministic visual specialists} (the OpenCV position-slot detector above; a
ResNet size head), and the VLM is kept to the coarse prefix and relation typing where it is reliable.
This is the neuro-symbolic interface in practice --- learn where the network is reliable, compute
deterministically where it is not --- and it is \emph{why} the realized path is the specialist, not
the VLM. (The one VLM-side gain still open is direction, $57\% \to$, via subtype-contrastive data
augmentation; the slowest loop, lowest priority.)

\begin{figure}[ht]
  \centering
  \includegraphics[width=\textwidth]{vlm_profile.png}
  \caption{Per-field re-evaluation of the fine-tuned VLM (G8) and its downstream consequence. The
  VLM saturates the coarse prefix (storey, class) and the relation-present signal, but recovers the
  discriminating fields (position-slot, size) at 0\%. The realized spatial-address specialist closes
  the filler subset from a 6.6\% floor to 67.6\%, while selective prediction raises the answered
  subset to 80.6\%.}
  \label{fig:vlm}
\end{figure}

\paragraph{(2) The confidence is calibratable --- the routing premise holds, and is not assumed.}
Before any routing we gate on calibration \citep{guo2017calibration}: a routing layer is only
legitimate if its confidence tracks correctness. On the held-out fillers the raw detector confidence
is positively discriminative (AUROC \textbf{0.80}: more-confident extractions are more often correct)
and only moderately mis-calibrated (ECE \textbf{0.206}); temperature scaling reduces it to
\textbf{0.172} (Figure~\ref{fig:caldiag}). The premise is established empirically, not asserted. We
note this gate is load-bearing: had the confidence been anti-correlated, no monotone recalibration
could rescue it, and the gain would have shifted to the schema-alignment and visual specialists (the
contingency we named in advance and did not need to invoke).

\begin{figure}[ht]
  \centering
  \includegraphics[width=0.85\textwidth]{calibration_diag.png}
  \caption{Calibration gate. The raw detector confidence is positively discriminative (AUROC 0.80)
  and moderately mis-calibrated (ECE 0.206); temperature scaling reduces ECE to 0.172. The reliability
  diagram shows the premise holds before any routing.}
  \label{fig:caldiag}
\end{figure}

\paragraph{(3) An honest negative: continuous reweighting is a no-op --- the calibration does not pay
off as a rerank weight.} With storey and class contributing integer agreement and the slot acting as
the \emph{finest} tiebreaker, any strictly-positive weight on a slot match induces the identical
ordering within a storey$\times$class bucket; the hard match (weight 1), the raw-confidence weight,
and the temperature-calibrated weight all yield the same \textbf{67.6\%} Top-1. Reweighting cannot
reorder what only a \emph{removal} of the term would change. We report this rather than bury it: the
value of the calibrated confidence is not in the soft weight.

\paragraph{(4) The payoff is selective prediction --- the practical mechanism and the triage value
proposition.} The calibrated confidence pays off as a \emph{deferral} threshold. Sweeping it traces a
coverage--accuracy curve (Figure~\ref{fig:rerank}): deferring the least-confident ${\sim}20\%$ of
cases lifts Top-1 on the answered subset from \textbf{67.6\% to 80.6\%}. The system does not return a
confident wrong GUID; it abstains and surfaces the candidates. A worked case makes the mechanism
concrete (Figure~\ref{fig:demo}): the detector predicts slot 1 of 10 where the truth is 8 of 10 ---
an error --- but its calibrated confidence is \textbf{0.05}, below threshold, so the case routes to
\emph{defer} rather than to a confident mistake. For a triage tool whose output carries downstream
consequence, ``here are the nine candidates'' at a known accuracy is the correct behavior, and the
coverage--accuracy curve --- not a single point estimate --- is the honest way to report it
\citep{geifman2017selective}.

\begin{figure}[ht]
  \centering
  \includegraphics[width=0.85\textwidth]{calibrate_rerank.png}
  \caption{Selective prediction. Sweeping the deferral threshold traces a coverage--accuracy curve;
  deferring the least-confident ${\sim}20\%$ lifts answered-set Top-1 from 67.6\% to 80.6\%.}
  \label{fig:rerank}
\end{figure}

\begin{figure}[ht]
  \centering
  \includegraphics[width=\textwidth]{case_AP_SK_092.png}
  \caption{A worked DEFER case. The detector predicts slot 1 of 10 where the truth is 8 of 10 --- an
  error --- but its calibrated confidence (0.05) is below threshold, so the case routes to defer and
  surfaces the candidate pool rather than emitting a confident wrong GUID.}
  \label{fig:demo}
\end{figure}

\section{A representational prerequisite: the address must be image-recoverable, not model-arbitrary}

Realizability has a precondition that the oracle analysis hides. The ground-truth slot index in the
source data is numbered along each wall's \textbf{IFC local-X axis} --- a modelling artefact whose
sign is invisible in any image. A detector reading the plan cannot recover which end of the wall the
modeller chose as ``first''; scored against that convention, the detector and ground truth disagree on
\textbf{16 of 35} fillers, \emph{all} exact mirrors $i \mapsto M-1-i$ (the count $M$ is identical).
The index, so defined, fails the address's own \emph{image-recoverable} criterion. The remedy is to
fix the slot's orientation by an \textbf{image-coordinate} convention --- number from a world
reference direction shared by both the detector and the ground truth --- under which the disagreement
vanishes and joint accuracy rises from an apparent 34\% to the true \textbf{74\%}. This is not
bookkeeping: it is a substantive constraint on what a deployable spatial address may be. An address
field is realizable only if it is defined in a frame the evidence carries; a field keyed to an
arbitrary modelling choice is, by construction, unrecoverable. The same property is what lets the
index survive the eventual patch-upload setting --- a north-up plan crop encodes the reference
orientation, whereas the IFC local axis never appears in any image.

\section{The honest boundary of the claim}

Three caveats bound the result. \emph{Statistical power:} the calibration and coverage--accuracy
estimates rest on 35 addressable fillers; the curve is clean down to ${\sim}0.74$ coverage and noisy
below, and the temperature estimate carries small-$n$ uncertainty --- we report the operating point
(defer ${\sim}20\% \to 80.6\%$), not a precise optimal threshold. \emph{Scope:} we realize and
calibrate the single highest-value extractor (the position-slot). The wall fingerprint is demonstrated
at oracle level only --- and its realization is not merely deferred but \textbf{blocked}: its
load-bearing fields (\texttt{connection\_degree}, \texttt{length\_band}) depend on IFC
wall-\emph{instance} boundaries the floorplan render does not encode (collinear instances merge into
one poch\'e; measured length-band recovery 5/17, realized Top-1 $\approx$ floor). This is the negative
case of the same image-recoverability constraint the position-slot satisfies --- the wall address is
oracle-discriminative but not image-realizable --- and it is why realization is scoped to the one
recoverable extractor; the remaining ``other''-class descriptors are genuine future work.
\emph{Mechanism honesty:} the headline is selective prediction, not soft reweighting --- the latter we
report as a measured no-op, and we explicitly do \emph{not} lead with ``we added calibration'', which
the calibratability gate supports but does not, on its own, deliver.

\section{Conclusion}

RQ2 asked how much of the oracle ceiling is realizable, and how a noisy address is used without losing
recall. The answers are concrete and, where they cut against the obvious design, honest. Recall is
preserved by making the address a soft prior inside a recall-fixed pool, never a hard filter (the
$\prod r$ collapse forbids the latter). One real extractor closes most of the oracle gap ($6.6 \to
67.6\%$ Top-1). Its confidence is calibratable (AUROC 0.80; ECE $0.206 \to 0.172$), so the routing
premise holds. Reweighting that confidence is a no-op; its genuine payoff is selective prediction,
where deferring the least-confident fifth lifts answered-set accuracy to 80.6\%. And realizability
itself imposes a representational constraint --- the address must be defined in an image-recoverable
frame --- without which the apparent accuracy is halved by a convention the camera cannot see. The
mechanism is therefore not ``deterministic retrieval'' but its disciplined replacement:
\textbf{auditable, deterministic execution over a calibrated, recall-safe, selectively-predicted
spatial address.}
